\section{抽样分布连接数据与推断}
\label{sec:11-Mathematical-Statistics/01-Statistical-Models/03}

一次 A/B 实验跑完两周，实验组转化率 \(5.4\%\)，对照组 \(5.1\%\)。会议室里的问题很短：这 \(0.3\) 个百分点，是新版本带来的，还是这两周恰好抽到了这批人？

两边各有一个点估计，上一节的图能解释它们为什么会不一样，却还不能判决。要判决，需要知道那团点摊得有多开：如果两组之间即使毫无差别，\(0.3\) 个百分点也经常出现，那这次的 \(0.3\) 说明不了什么；如果在无差别的前提下这么大的差距很少发生，结论就是另一个方向。回答这个``经常还是很少''，靠的是估计量的抽样分布。

麻烦在于，实验只做了一次。``经常出现''说的是那些没有发生的重复，而我们手上只有发生了的这一次。

\begin{center}
\begin{tikzpicture}[x=1.25cm,y=0.95cm,font=\small,>=stealth]
  \draw[->,black!55] (-2.0,0) -- (2.3,0) node[right,font=\footnotesize] {\(\hat\theta\)};
  \fill[red!70!black] (0.6,0) circle (3.0pt);
  \node[font=\footnotesize,red!70!black,anchor=south] at (0.6,0.18) {\(\hat\theta_{\mathrm{obs}}\)};
  \node[font=\footnotesize,align=center,black!75] at (0.1,-0.95) {你拥有的：一次实验，一个数};
  \begin{scope}[shift={(6.3,0)}]
    \draw[->,black!55] (-2.6,0) -- (2.9,0) node[right,font=\footnotesize] {\(\hat\theta\)};
    \draw[blue!60!black,very thick] plot[domain=-2.5:2.5,samples=80] (\x,{2.0*exp(-\x*\x/2)});
    \draw[dashed,black!70] (0,0) -- (0,2.15);
    \node[font=\footnotesize,black!75,anchor=south] at (0,2.18) {\(\theta\)};
    \foreach \px in {-1.7,-1.05,-0.55,0.15,0.95,1.45}
      \draw[black!45] (\px,0.16) circle (2.0pt);
    \fill[red!70!black] (0.6,0) circle (3.0pt);
    \draw[<->,black!55] (-1.0,1.25) -- (1.0,1.25);
    \node[font=\footnotesize,black!70,anchor=south] at (0,1.3) {一倍标准误};
    \node[font=\footnotesize,align=center,black!75] at (0.15,-0.95) {抽样分布：重复这次实验\\会得到的那些数};
  \end{scope}
\end{tikzpicture}

\emph{左边是发生过的事，右边那条曲线由没有发生过的重复定义；把左边的点放进右边，才谈得上``这个差异是不是噪声''。}
\end{center}

图里的空心小圈一个也没有发生。它们是同一个生成机制、同一个样本量下本来可能抽到的结果，是平行世界里的报表。抽样分布就是这些平行世界的分布，而整套频率学派推断——标准误、置信区间、\(p\) 值、检验阈值——全部是在这条曲线上量出来的。一次实验只交给你一个红点；把红点放到曲线上，它才从一个数字变成一句带不确定性的断言。

\subsection{抽样分布是一个思想实验}

先把定义摆清楚。估计量 \(\hat\theta=T(X_1,\ldots,X_n)\) 在重复抽样下的概率分布，称为它的抽样分布。它跟另外两个分布容易混：\(X_i\) 自身的分布 \(P_\theta\) 说的是单条记录长什么样，本次样本的经验分布说的是这 \(n\) 条记录实际取了哪些值，而抽样分布说的是一个汇总量在不同批次之间怎么变。三者的自变量不同，形状也毫无关系——\(X_i\) 可以是 0/1 两点分布，\(\bar X\) 的抽样分布却接近连续的钟形。

这条曲线依赖于未知参数，写全了是 \(\hat\theta\) 在 \(P_\theta\) 下的分布，随 \(\theta\) 移动而移动。所以严格说我们从来没有一条抽样分布，只有一族；后面所有的技术动作，要么是在这一族里找一个不依赖 \(\theta\) 的量（枢轴量），要么是把 \(\theta\) 换成估计值（近似）。

它还依赖于``重复''二字的具体含义，这一点常被跳过。重复实验时什么跟着变、什么固定不变，需要事先说明。回归里的特征矩阵就是典型：把它看作随机抽取，抽样分布里就平均掉了设计本身的随机性；条件于已经观察到的设计矩阵，抽样分布只反映响应噪声。剂量由实验者固定时条件视角更自然，要把结论推广到新的用户特征分布时无条件视角更合适。两种做法的标准误目标不同，谁也不比谁更正确，但报告时必须说清楚重复的是哪一件事。

\subsection{两个模型里它可以精确算出来}

上一节的评测例子里，判对与否是 \(X_i\overset{\iid}{\sim}\Bernoulli(\theta)\)，判对条数 \(S=\sum_i X_i\) 的分布由定义直接给出，即 \(n\) 次独立试验的成功次数，

\[
S\sim\Binomial(n,\theta),\qquad\hat\theta=\frac Sn.
\]

这就是那条曲线的精确形状，离散、取值在 \(\set{0,1/n,\ldots,1}\) 上。由 \(\E_\theta[S]=n\theta\) 与 \(\Var_\theta(S)=n\theta(1-\theta)\)，两边分别除以 \(n\) 和 \(n^2\)，

\[
\E_\theta[\hat\theta]=\theta,
\qquad
\Var_\theta(\hat\theta)=\frac{\theta(1-\theta)}n,
\qquad
\SE_\theta(\hat\theta)=\sqrt{\frac{\theta(1-\theta)}n}.
\]

下标 \(\theta\) 不是装饰，它提醒这三个量都随未知参数变化；实践中代入 \(\hat\theta\) 得到估计标准误，于是又叠了一层近似。

正态模型给出另一个精确解。\(X_i\overset{\iid}{\sim}N(\mu,\sigma^2)\) 时独立正态的线性组合仍是正态，于是

\[
\bar X\sim N\!\left(\mu,\frac{\sigma^2}n\right),
\qquad
Z=\frac{\bar X-\mu}{\sigma/\sqrt n}\sim N(0,1).
\]

\(Z\) 的分布不含任何未知参数，这正是它能用来构造精确区间的原因。把 \(\sigma\) 换成样本标准差 \(S\) 之后，比值的分母也随机了，分布随之改变，在正态模型下变成自由度 \(n-1\) 的 Student \(t\)。这不是一条需要背诵的小样本修正口诀，它是估计干扰参数所付的利息：\(\sigma\) 未知这件事没有消失，只是从区间的中心挪到了它的形状上。

条件对一遍。Binomial 那一行用到独立与同分布，两者缺一，方差公式就不成立；正态那一行还额外用到 \(X_i\) 精确正态，这一条在有限样本下几乎从不成立。下面这个近似就是用来替换最后这个条件的。

\subsection{中心极限定理把桥修到大多数模型上}

大部分实际数据既不是 Bernoulli 也不是正态。中心极限定理在方差有限的 \(\iid\) 前提下给出

\[
\frac{\sqrt n(\bar X-\mu)}{\sigma}\overset{d}{\longrightarrow}N(0,1),
\]

于是那条曲线在大样本下总可以近似当成正态来量，无论单条记录本身的分布多难看。这条定理是整个应用统计的承重墙，推导见\hyperref[sec:05-Probability/08-Limit-Theorems/03]{``中心极限定理与 Gaussian 近似''}。

但``\(n\ge30\) 就可以用正态''不是定理，只是经验口诀。收敛速度取决于偏度、尾部厚度、离群点的多少，也取决于你关心分布的哪个部位：均值附近收敛得快，而 \(99.9\%\) 分位数那种极端区域可能需要大得多的样本量。转化率这类稀有事件尤其要当心，\(np\) 太小时二项分布明显偏斜，正态近似给出的区间会越过 0 或者失衡。方差无限时经典中心极限定理直接失效，重尾数据的样本均值本身就不稳定；数据相关时 \(\sigma/\sqrt n\) 也不再对，需要时间序列或聚类版本的标准误。样本量大并不自动买到正态近似。

\subsection{标准误只覆盖抽样这一种误差}

标准差描述单条记录 \(X_i\) 的分散，标准误描述汇总量 \(\hat\theta\) 的分散，两者量纲相同而含义相距很远。样本均值的标准误按 \(1/\sqrt n\) 收缩，个体之间的差异却不会因为你多收了数据而缩小一分。多收数据让均值的位置更确定，不让世界变得更整齐。

更要紧的是这个数没有覆盖的部分。抽样分布是在假定的生成机制内部计算出来的，它度量的只是``同一个机制重复一次会有多少晃动''。抽样框有偏、模型错设、埋点漂移、目标人群改变——这些在\hyperref[sec:11-Mathematical-Statistics/01-Statistical-Models/01]{``数据来自怎样的生成机制''}里讨论过的问题，一律不在标准误里。它们让整条曲线整体挪位，而曲线的宽度对此一无所知。样本量做到千万级时，标准误可以小到读数的第四位，此时误差几乎全部来自它没覆盖的那一部分。

\subsection{把 0.3 个百分点放进抽样分布}

回到开头。两组独立，各 \(n\) 人，转化率差 \(D=\hat p_1-\hat p_0\)，独立性让方差可加，

\[
\SE(D)=\sqrt{\frac{p_1(1-p_1)}{n}+\frac{p_0(1-p_0)}{n}}.
\]

每组 5000 人时，代入观测比例得 \(\SE(D)\approx0.45\) 个百分点。观测到的 \(0.3\) 还不到一倍标准误，落在曲线的正中央那一带——即使两个版本毫无差别，这么大的差距也会经常出现，实验没有给出可以行动的信息。每组 40000 人时同一个公式给出 \(\SE(D)\approx0.16\) 个百分点，同样的 \(0.3\) 变成了将近两倍标准误，落到了曲线的边缘。

同一个 \(0.3\) 个百分点，两种读法。差别不在效应本身，在那条曲线的宽度：宽的时候它是噪声的常态，窄的时候它是异常。这也解释了实验平台为什么坚持先算样本量再开实验——样本量决定曲线宽度，而曲线宽度决定了多大的效应才有可能被看见。把``落在曲线的哪个位置''翻译成一个可以照做的判据，是\hyperref[sec:11-Mathematical-Statistics/05-Hypothesis-Testing/01]{``两类错误与功效共同定义检验''}和\hyperref[sec:11-Mathematical-Statistics/04-Interval-Estimation/01]{``置信度属于程序，而非这次参数''}要处理的事情；这里只需要记住，被翻译的那个东西就是抽样分布。

\subsection{没有公式时就把思想实验模拟出来}

上面两个模型能写出闭式解，是因为它们简单。换成中位数、分位数、AUC、两个模型指标之比、一整条特征工程加训练加评测的流水线，抽样分布通常写不出来。

模型能模拟时有一条出路：固定一个候选参数，反复生成样本、反复计算统计量，用这些结果的经验分布近似那条曲线。这对核对推导、看清有限样本下的偏斜特别有用。代价是模拟的输入就是你假设的模型——它能告诉你``若世界如此，统计量会怎样波动''，不能反过来证明世界如此。

参数未知时还有一条更省事的出路，就是 bootstrap 的动机。既然重复实验意味着从 \(P\) 里重新抽 \(n\) 个观测，而 \(P\) 未知，那就用手上这份数据的经验分布顶替 \(P\)：从原样本里有放回地抽 \(n\) 条，重算统计量，重复几千次。得到的分布近似的正是图中那条曲线。整节反复强调抽样分布是想象出来的，bootstrap 做的事情是把这个想象在数据上重放一遍——平行世界仍然没有发生，但可以用重采样造出一批替身。它依赖经验分布确实接近 \(P\)，因而在小样本、重尾、极端分位数上照样会失手，完整讨论见\hyperref[sec:11-Mathematical-Statistics/06-Resampling-and-Model-Selection/01]{``Bootstrap 近似估计量的抽样波动''}。

这一单元到此把地基铺完了：机制的声明决定结论替谁说话，参数与统计量分处两类对象，抽样分布把一次观测接到一句带不确定性的断言上。接下来的问题变成造这个断言的工具本身——同一个 \(\theta\) 可以有许多估计量，哪个更好，好在哪里，先从\hyperref[sec:11-Mathematical-Statistics/02-Point-Estimation/01]{``偏差、方差与均方误差''}开始拆。
